\chapter{Marco teórico}\label{cap4}

\section{Diseño y tipos de biorreactores}

Un biorreactor es un dispositivo capaz de proporcionar un entorno homogéneo y
favorable para que se lleven a cabo las reacciones químicas que fomentan que los
organismos tratados produzcan los metabolitos deseados. El diseño de un
biorreactor involucra contemplar parámetros como concentración, tasa de flujo
volumétrico y volumen. Además, es necesario considerar el uso de dispositivos de
medición, como sensores de pH u oxígeno, puesto que diferentes variables juegan
un papel importante en el crecimiento del organismo \cite{Jaibiba2020}.

Entre los diferentes tipos, los biorreactores tipo airlift, se caracterizan por utilizar la
turbulencia creada por gas para mezclar el contenido. El aire se introduce desde el
fondo del reactor, lo que resulta en un movimiento circulatorio del contenido dentro
del reactor que ayuda a obtener una mayor transferencia de oxígeno.
Además, presentan ventajas como bajo consumo energético, altas tasas de
transferencia de calor y masa, menor estrés por cizallamiento debido a la ausencia
de agitadores mecánicos, y costos de construcción y operación reducidos \cite{Esperanca2023}.

\section{Factores clave en el uso de biorreactores airlift}

La hidrodinámica representa uno de los factores críticos en el diseño de
biorreactores airlift, puesto que afecta en el crecimiento del organismo y la viabilidad
del proyecto. Suelen evaluarse parámetros hidrodinámicos como la retención del
gas ($\varepsilon$), la velocidad del líquido ($V_L$) y la tasa de corte ($\dot{\gamma}$). En la actualidad se utiliza
la dinámica de fluidos computacional (CFD) para estos análisis, puesto que
significan un ahorro de tiempo, energía y reactivos \cite{Esperanca2023}.

El cultivo de hongos filamentosos representa desafíos en lo respectivo al diseño del
biorreactor. La reología del medio, transferencia de oxígeno y estrés por
cizallamiento son factores por considerar relacionados con la eficiencia de su
cultivo. También es necesario considerar la morfología del hongo filamentoso,
puesto que la selección del tipo de reactor, en específico los mecanismos para la
mezcla, dependen de esta \cite{Gomes2023}.

En particular, los ascomicetos coprófilos requieren de condiciones específicas de
pH que deben variar según su etapa de desarrollo, así como una humedad y
oxigenación elevadas. Además, se ha encontrado también que el desarrollo de sus
ascomas es estimulado por la luz. Las características químicas del medio de cultivo
tienen un impacto en el desarrollo del hongo, por lo que es necesario asegurar la
dispersión homogénea de los nutrientes \cite{Bills2013,Gomes2023}.

\section{Escalamiento y automatización de biorreactores}
Para el diseño geométrico de biorreactores airlift a escala piloto, la literatura
reporta relaciones estandarizadas entre las dimensiones principales del sistema.
La relación entre el diámetro del tubo de tiro ($D_B$) y el diámetro interno del
fermentador ($D_F$) se sitúa típicamente entre 0.5 y 0.6, mientras que la relación
entre la altura del fermentador ($H_F$) y su diámetro puede variar entre 3 y 15
dependiendo de la escala y aplicación \cite{Moresi2025}. 
Otro parámetro por considerar es el mezclado y transferencia de masa, puesto que
el tamaño seleccionado debe garantizar una transferencia eficiente de nutrientes y
oxigeno sin generar esfuerzo cortante excesivo. Por otra parte, biorreactores de
mayor capacidad implican costos iniciales y de operación más elevados, además
de requerir procedimientos de esterilización más complejos y costosos \cite{Moresi2025,Micheletti2006}.

La automatización mediante dispositivos eléctricos, como placas de desarrollo y
sensores (temperatura, presión, etanol, luminosidad y oxigenación) permite
monitorear y controlar las condiciones internas del sistema. Como consecuencia,
se incrementa la productividad del sistema y se asegura una calidad constante en
el producto generado \cite{Paredes2024}.
